\section{The forward marker theorem}\label{sec:forward}

This section concerns pure markers.  Its target is a product-interval
statement forcing a bounded prefix and many actual point stabilizers
in an ordinary site group.  The individual group-type arguments that
follow are exclusions under the assumption that this point configuration
is absent; their final contradiction again uses the protected path
bound.  The normal and mixed-projection tests below will also apply
internally when the simple sites change.

\begin{definition}[Pure marker lattice]\label{r:pure-markers}
For a finite set $T$ of $t$ markers, let $F_T$ have labels
\[
 (I,V),\qquad I\subseteq T,\quad V\le S_{T\setminus I}.
\]
Choose a set $Z$ disjoint from $T$, with $|Z|=t+1$ or larger, and
interpret this label as the subgroup
\[
 S_{Z\cup I}\times V\ \le S_{Z\cup T}.
\]
Order, meet, and join are subgroup containment, intersection, and
generation in this model.  The markers in $I$ are called absorbed.
We write
\[
 X_I=(I,1),\qquad x_i=X_{\{i\}},\qquad
 M_D=X_{T\setminus D},\qquad
 C_D=(T\setminus D,S_D)\quad(D\ne\varnothing).
\]
The bottom and top are $X_\varnothing$ and $X_T$.
We also write $F_t$ when only the cardinality matters.
\end{definition}

\begin{lemma}[Elementary marker properties]\label{r:marker-elementary}
The labels in \Cref{r:pure-markers} are closed under meet and join.
The coatoms are exactly the $C_D$, $\varnothing\ne D\subseteq T$.
The axis filter $[X_I,1]$ is $F_{T\setminus I}$, and the axis ideal
$[0,X_I]$ is $F_I$.  Moreover
\[
 \operatorname{ht}(F_t)=O\bigl(t\log\log(t+3)\bigr).
\]
\end{lemma}

\begin{proof}
The independent symmetric group on the absorbed part makes the
intersection again a group with independent full symmetric action
on its new absorbed part.  In a join, that part expands exactly
along the generated orbits meeting it: symmetric transpositions
there and their conjugates give the full symmetric group on the
expanded part, independently of the induced group on the remaining
markers.  Thus both operations stay within the indicated labels.
These rules depend only on the marker labels $(I,V)$, so the resulting
lattice is independent of the choice of $Z$ with $|Z|>t$.

If a proper label has unabsorbed set $D$, it lies below $C_D$.
The latter is the full stabilizer of the partition into $Z\cup(T\setminus D)$
and $D$.  These parts have unequal sizes, the first larger than
$t\ge|D|$; adjoining any element failing to preserve them generates
the full symmetric group by the same absorption argument.  Hence
$C_D$ is a coatom, proving the classification.  The descriptions
of the axis filter and ideal follow directly by treating $Z\cup I$
as the fixed absorbed set, or by restricting to $Z\cup I$,
respectively.

For the height use $|Z|=t+1$, embedding $F_t$ in the subgroup lattice
of $S_{2t+1}$.  Every strict step in a subgroup chain consumes at
least one prime factor, with multiplicity, of $(2t+1)!$.
Writing $n=2t+1$, the total number of such factors is
\[
 \sum_{p\le n}\sum_{a\ge1}\left\lfloor\frac{n}{p^a}\right\rfloor
 \le n\sum_{p\le n}\frac1{p-1}
 =O\bigl(n\log\log(n+2)\bigr),
\]
by the reciprocal-prime bound (the difference between $1/(p-1)$
and $1/p$ is summable).  This proves the estimate.
\end{proof}

\begin{lemma}[Marker normal test]\label{r:marker-normal}
For sufficiently large $t$, every exact finite group interval
isomorphic to $F_T$ or $F_T^{\mathrm{op}}$ has property~\((*)\).
\end{lemma}

\begin{proof}
Let $\omega=(I,V)$ label the base times a normal subgroup.
As in \eqref{r:left-modular}, it and all its projections to
subintervals satisfy the left-modular law in either orientation.
Indeed in actual subgroup order, with base $H_0$, normal
$N\trianglelefteq G_0$, and $H_0\le D\le E\le G_0$, the identity is
\[
 (DN)\cap E=D(N\cap E)=D\vee((H_0N)\cap E).
\]
Projection to $[D,E]$ gives $D(N\cap E)$, with
$N\cap E\trianglelefteq E$.  Self-duality of the left-modular
identity gives the reversed statement as well.  In particular the configuration
\begin{equation}\label{r:marker-forbidden}
 \omega\wedge b\le a<b\le\omega\vee a
\end{equation}
is impossible, also after any projection.

Suppose $I$ is nonempty and proper, with $|I|\ge2$.
Choose four positions $P_4$ in cyclic order such that membership
in $I$ is nonconstant, occurs in both parity classes, and is not
invariant under rotation by one.  Such a choice is possible for
large $t$: two chosen positions of $I$ can be adjacent, and at
least one of the other positions lies outside $I$; if the
complement has only one point, use three from $I$.
For the rotation $s$ put
\[
 a=(T\setminus P_4,\langle s^2\rangle),\qquad
 b=(T\setminus P_4,\langle s\rangle).
\]
The $I$-membership condition forces $\omega\wedge b\le a$.
Each parity orbit of $s^2$ meets the absorbed part of $\omega$,
so $\omega\vee a=1$.  Thus \eqref{r:marker-forbidden} holds.

If $I=\{j\}$ and $V$ moves some marker $k$, choose $P_4$ so that
$k$ has opposite parity from $j$ and its image under some element
of $V$ lies outside $P_4$.  That image is already absorbed in $a$;
the two parity classes are therefore absorbed in $\omega\vee a$.
Again the same rotation gives \eqref{r:marker-forbidden}.
If $V=1$, use three positions containing $j$, with $C_3<S_3$
in place of the two rotation groups and absorb the outside.
The join with the $C_3$ absorbs all three positions, while the
meet with the $S_3$ lies in $C_3$.  Thus nonempty proper $I$ is
impossible in all cases.

Now let $I=\varnothing$ and suppose $V$ moves $j$.
Project above $x_j$.  In the marker description of that axis
filter, $\omega\vee x_j$ has nonempty absorbed set because the
$V$-orbit of $j$ contains another point.  The preceding test makes
it the top.  Consequently $V$ is transitive on $T$.
Every marker atom $x_j$ complements $\omega$.

First take the forward orientation.  Complementation of a
base-normal multiple gives the interval isomorphisms
\[
 [0,x_j]\longrightarrow[\omega,1],\qquad u\longmapsto u\vee\omega.
\]
As $x_j$ is an atom, $\omega$ is a coatom; since its absorbed set
is empty, \Cref{r:marker-elementary} gives $\omega=C_T$.
Factor by the base core.  Choose a minimal missing normal subgroup
inside the putative normal subgroup; it still has a nonbottom
proper normal-multiple label, and the same argument identifies
that label with $C_T$.  Write this minimal normal as $N$ and the
base as $H$.  At a transposition label
\[
 0<d=(\varnothing,\langle(i\ k)\rangle)<C_T,
\]
the intersection $N\cap K_d$ is a nontrivial proper subgroup of
$N$, and $H(N\cap K_d)=K_d$.
For every $j$ the marker operations give
\[
 C_T\wedge(d\vee x_j)=d.
\]
Because $N$ is normal in the whole upper endpoint and lies in
$K_{C_T}$, this equality shows that $K_{x_j}$ normalizes
$N\cap K_d$.  The groups $K_{x_j}$ generate the upper endpoint.
Thus $N\cap K_d$ is a nontrivial proper normal subgroup of $N$
normal in that endpoint, contradicting minimality.

In the reverse orientation the normal projection instead yields
isomorphisms, in the displayed $F_T$ order,
\[
 [x_j,1]\longrightarrow[0,\omega],\qquad u\longmapsto u\wedge\omega.
\]
For $k\ne j$ the label $X_{\{j,k\}}$ is strictly above $x_j$;
therefore $X_{\{j,k\}}\wedge\omega>0$.
This meet is nontrivial only if $V$ contains the transposition
$(j\ k)$.  Hence $V=S_T$ and $\omega=C_T$.
But a supported three-cycle group containing $j$ cannot be in
the image of the indicated isomorphism.  Any label above $x_j$
whose meet with $C_T$ contains that cycle must absorb all three
positions; its meet then contains the full $S_3$, rather than
just the cycle group.  This contradicts surjectivity.

We have proved the normal endpoint assertion.  The lattice $F_T$
is nonmodular for large $t$: it contains the full subgroup lattice
of $S_T$, or one may use the rotation obstruction just exhibited.
In the faithful base-coset action after factoring the core, every
minimal normal subgroup is transitive because it supplements the
base.  An abelian minimal normal identifies the interval with an
invariant-subspace lattice and makes it modular.  Two distinct
nonabelian minimal normals commute; supplementation by one makes
the base induce all inner actions on the other, so intersections
there are normal subproducts of simple factors, again forming a
modular lattice.  These contradictions give the unique nonabelian
monolith required by property~\((*)\).
\end{proof}

\begin{lemma}[Exclusion of a mixed marker base]\label{r:marker-mixed}
For sufficiently large $t$, suppose that
\[
 F_T\cong[P,S^b]^\Delta
\]
in forward order, with $S$ nonabelian simple and $\Delta$
transitive on its factors.  If $c$ denotes coordinate-product
closure, then
\[
 P<c(P)<S^b
\]
is impossible.
\end{lemma}

\begin{proof}
Assume the mixed case and let $\omega=(I,V)$ label $c(P)$.
The projections of $P$ are nontrivial and proper at every site,
by factor transitivity.  A label joining with $\omega$ to the top
must be subdirect: adjoining $c(P)$ does not change its site
projections.  For two distinct comparable subdirect labels, their
intersections with $\omega$ are distinct.  Indeed their tied
coordinates and restricted twisting maps are realized using the
nontrivial full projections of $P$ to the factors of $c(P)$;
a strict refinement of ties gives a strict change of the
intersection.  Conversely, comparable subdirect groups with
unchanged ties and twists coincide.

Apply the rotations in the proof of \Cref{r:marker-normal}.
They give $\omega\wedge b\le a<b$ with both $a$ and $b$ joining
$\omega$ to the top.  Hence both are subdirect, whereas
\[
 a\wedge\omega=b\wedge\omega,
\]
contradicting the preceding strictness.  This eliminates nonempty
proper $I$, including the singleton cases of that proof.
The top absorbed set is impossible because $\omega$ is proper.

If $I=\varnothing$ and $V$ moves $j$, consider $c(x_j)$.
It contains $\omega\vee x_j$.  In the axis filter above $x_j$,
apply the nonempty-absorption test just used if the product
closure is still a proper mixed label.  If $x_j$ itself is already
a product, its closure equals it; this cannot happen because it
contains $\omega\vee x_j>x_j$.  Thus $c(x_j)$ must be the top.
So $x_j$ is subdirect.  Its join with $\omega$ is then the top:
$\omega$ has independent nontrivial supports, whose conjugates
under a subdirect group generate each simple site by simplicity.
In parameter terms this makes $V$ transitive.  All $x_j$ are
therefore subdirect and have $x_j\wedge\omega=0$.
Strictness for the comparable subdirect labels
$x_j<X_{\{j,k\}}$ gives
$X_{\{j,k\}}\wedge\omega>0$; as before, $V$ contains every
transposition, and so $V=S_T$.

Each meet $x_j\wedge\omega=0$ is the same subgroup $P$.
Write $c(P)=\prod_o T_o$ over the simple sites, where every $T_o$ is
nontrivial.  On a tied pair of sites, the two-coordinate projection
of $P$ is the graph of the identifying isomorphism restricted to
$T_o$; on an untied pair it is the whole product.  Thus $P$ determines
the ties and the restrictions of their twisting maps to the $T_o$,
although it need not determine the full maps on $S$.  Every invariant
overgroup of a subdirect label is obtained by an operator-invariant
refinement of its ties, retaining the twists on the remaining parts.
Its intersection with $c(P)$ depends only on that refinement and the
restricted twists.  Consequently meeting with $\omega$ has the same
image on every filter above $x_j$.

This is false for $\omega=C_T$.  Fix a triple of markers.  If $j$
lies outside it, joining its supported three-cycle label with $x_j$
and then meeting $C_T$ recovers exactly that three-cycle label.
If $j$ lies inside the triple, any label above $x_j$ whose meet with
$C_T$ contains the cycle must absorb all three positions; its meet
then contains the full $S_3$.  It cannot yield just the three-cycle
subgroup.  A triple and an outside marker give the contradiction.
\end{proof}

\begin{lemma}[Ordinary site interval for product markers]
\label{r:marker-site}
Suppose $F_T=[P,S^b]^\Delta$ exactly in forward order, with
$\Delta$ transitive on the simple factors and $P$ a coordinate
product satisfying $P\ne1$ or $b=1$.  For sufficiently large $t$,
all coatoms and all axis labels are products.  At one site, let
$Q$ be the projection of $P$ and let
$\Lambda\le\operatorname{Aut}(S)$ be the image of the site
stabilizer.  Then $\Lambda\cap S\le Q$, and the product labels
form the ordinary interval
\[
 [H,S\Lambda],\qquad H=\Lambda Q.
\]
This correspondence preserves meets and joins in the full invariant
interval.  We denote its subgroup at a product label $u$ by $K_u$.
\end{lemma}

\begin{proof}
A proper subgroup projecting fully at every simple site cannot
contain the independent nontrivial base projections: their
conjugates normally generate the full simple sites.  At $b=1$
this is automatic.  Hence each coatom has proper product closure
and equals that closure by maximality.
For a proper coatom let $T_1$ be its projection at the fixed site.
It is nontrivial and proper.  The group
$T_1(\Lambda\cap S)$ is still invariant and transports to every
site.  It is proper, since equality to $S$ would imply
$T_1\trianglelefteq S$.  Maximality consequently gives
$T_1(\Lambda\cap S)=T_1$.
The singleton-omission coatoms $C_{\{i\}}$ intersect to the base,
so their site projections give $\Lambda\cap S\le Q$.
The same intersections over $i\notin I$ give $X_I$; thus every
axis label is a product.

An invariant site projection containing $Q$ transports uniquely
to a $\Delta$-invariant coordinate product.  As
$\Lambda\cap S\le Q$, these projections are exactly the
intersections with $S$ of the overgroups of $H$ in $S\Lambda$.
An overgroup is recovered by multiplying its intersection with
$S$ by $H$.  Finally, coordinate products are closed under
intersection and generation in the original simple power,
proving the asserted compatibility.
\end{proof}

\begin{theorem}[Forward marker theorem]\label{r:forward-marker}
For every nonnegative integer $m$, there are integers $a(m)$ and
$t_0(m)$ such that the following holds whenever $|T|\ge t_0(m)$.
Every exact forward realization
\[
 F_T=[P,S^b]^\Delta
\]
with $S$ nonabelian simple, $\Delta$ transitive on the factors,
and invariant coordinate product $P\ne1$ (or $b=1$), has the
following configuration in the ordinary site interval of
\Cref{r:marker-site}.  There is $A\subset T$ with $|A|\le a(m)$
such that $K_{M_A}$ has an actual high-degree alternating simple
normal subgroup $B_A\le S$ satisfying
\[
 HB_A=K_{M_A},\qquad
 C_{K_{M_A}}(B_A)\le H,
 \qquad C_{K_{M_A}}(B_A)\text{ is soluble}.
\]
In the almost simple quotient
$K_{M_A}/C_{K_{M_A}}(B_A)$, at least $m$ distinct labels
$M_{A\cup\{i\}}$, with fresh indices $i\in T\setminus A$,
are full stabilizers of individual points in the natural
alternating action.  The constants $a(m)$ and $t_0(m)$ are independent
of $S,b,\Delta,P$ and the realization; the high-degree threshold is
the fixed one used in the structural exclusions.
\end{theorem}

The proof of \Cref{r:forward-marker} occupies the rest of this
section and \Cref{sec:forward-completion}.  Its initial site
reduction is \Cref{r:marker-site}.  The normal test and mixed-base
exclusion apply whenever new simple sites occur internally.
The bounded-type step uses exactly the sufficiently-general
bounded-complexity input established in \Cref{sec:bounded}; the
high-type steps retain the protected root and tensor conditions
of \Cref{sec:weights}.
